\section*{AI usage and my personal words}
This paper was written with ``assistance'' from (in fact, entirely by)
ChatGPT-5.6-Sol, ChatGPT-6-Astra, and ChatGPT-6.1-Sol.

As of October 10, 2026, I now claim the correctness of the paper because
the most important theorem compiles without \texttt{sorry} or extra axioms.
See the repository for details. The remaining risks are:
\begin{enumerate}
\item I have only falsified the rigidity conjecture; I have not done the
other propositions, such as the Cantor family.
\item The translation of the theorem might not be entirely risk-free.
AI has checked it, and it is highly unlikely to mislead us, but one can
never be certain.
\item Mistranslations might occur during the process, and the manuscript's
lemmas and exact proof details might not be entirely correct. The
formalization guarantees the existence of the band and the counterexample,
but the other steps might not be accurately represented. The manuscript
will undergo many edits after the formalization, and these edits might
deviate further from the formalization's route. In fact, this faces the
same problem as the Navier--Stokes equation proof. Yet, if we treat
natural-language manuscripts as sketches and roadmaps, guiding a correct
Lean proof already guarantees their trustworthiness. You can always ask
an AI to fill a gap if you find one; having a correct roadmap gives you
confidence that a way to fill it can be found.
\item The Lean kernel might contain bugs. I do not think this is likely
to be a problem, and if there really is a problem, it is hardly mine
alone. If the sky falls, the tall guys will hold it up.
\end{enumerate}

I had not noticed this problem until Levent Alpoge produced a
counterexample to the Carath\'eodory conjecture, followed almost immediately
by YTD and $S^6$. At that point, I began to believe that AI could solve
problems in differential geometry. I began asking GPT which problem on
Ghomi's list of open problems could be next (Carath\'eodory was on the
list). GPT suggested this one because the analytic case is true, the
polyhedral case is false, and the smooth case is unsolved. This resembled
the Carath\'eodory conjecture, where the analytic case is also true but the
smooth case was ultimately shown to be false. I then asked chat mode to
solve the problem. It found a counterexample to Problem~1.1 in Ghomi's
2025 paper, showing that the simple route toward proving rigidity was
false. It did not give any other promising result. Several days later,
I decided to make an all-out effort and downloaded Danus, running it for
28 hours and using about 140\% of the weekly allowance on the \$200 Pro
plan. Then Astra came out, and it took two hours and 25\% of the Pro
allowance to find the counterexample. Although the counterexample has
changed considerably, its idea has survived into the final result.
Evidence of this timeline can be found in the
\href{https://github.com/zhangyuetong/smooth-tight-nonrigidity/blob/6a0758c99077e2c7fd31c3de6ee3452b3f6df8d3/timestamp/README.md}{timestamp archive}.

I was shocked, so I asked it to write out the result and considered
giving it to my mentor or publishing it quickly. Yet the paper was not
a clean formula like the Carath\'eodory result: it was over 300 pages
long and practically impossible to read. The main construction and
proof involved PDEs, and I had not even taken a course in PDEs as an
undergraduate. For the next several weeks, I worked continuously to
make the paper simpler, more readable, and more rigorous:
\begin{itemize}
\item Around this time, I also learned that the mathematics community
advocates building mathematical theories and regards open problems as
golden geese. I therefore tried asking AI to ``generalize the theorems
introduced into lemmas and `theories'.'' The point was to make AI find
the ``essence'' of the new theorems, simplify and generalize their
conditions, and identify where they share the same proof. This could
have four outcomes: the theories might actually be useful; the paper
might become more reusable because the ``theories'' would remain even
if part of the construction failed; the theorems might become more
trustworthy because fewer conditions could reduce the cognitive burden
on both AI and humans; and, if AI actually built theories from scratch,
it might think more logically and reduce errors. Yet this approach
essentially failed. Most of the construction was too specific, and
the work did not produce useful theories. It also expanded the paper
to 400 pages.
\item I read a story online about someone who made substantial progress
and had various AI systems check it, but the AI overlooked a fatal gap,
leaving a month's effort in vain. This made me think that I might need
to formalize my paper. Yet I did not have enough tokens, and the
infrastructure for differential geometry was inadequate. The AI made
no progress with the formalization. This was the situation in September.
After OpenAI released the new infrastructure, I succeeded.
\item The final successful route was consulting ChatGPT Pro in chat
mode. This led to significant simplification, rewriting, and checking
of the paper. First, I asked it to rewrite the paper to remove the
branches, which brought it to 200 pages. Then I asked it to brainstorm
simplifications, which brought it to 100 pages. Next, I asked it to
brainstorm more intensively and find a good route, and then let it
choose one; this brought the paper to 16 pages. Finally, I asked it
to fill in the gaps, which brought it to around 30 pages. Its length
has not changed much since then.
\end{itemize}

For now, I do not know how I should take credit. In fact, I did not
do anything---or perhaps I have become one of those black-hearted
mentors exposed online who suppress their students and take all the
credit. I am an undergraduate who could not even understand the paper's
basic tools, yet buying a \$200 GPT subscription gave me a paper that
solved a long-standing open problem, one I might never have been able
to reach in a lifetime without AI. Previously, this result might have
earned me a paper in JDG, but now I do not even think anyone cares:
more important problems are being solved every day, and by the time
I wrote this, OpenAI had released over 700 results more important than
this one.

Yet, as an undergraduate without an original paper (and perhaps without
one ever, since we do not know what AI will do next year), I care about
receiving credit for at least doing something. I also care about the
effort I put into this paper, although it was not much. This feels
ironic because my dilemma somehow has the same structure as mathematics
itself. For the reasons above, every day that I delay publishing the
paper is another day when someone else might obtain the result with
a stronger model, leaving my month's effort in vain. In fact, Ding's
paper of October 6 might be a sign that he has been attacking this
problem with AI as well. Without me, he could use the same AI to
publish this result within a month, or even within days. Yet, for
a mathematician like Ghomi, a year of effort on this problem since
2025 might also be rendered vain by the publication of my paper.

The accompanying GitHub repository contains the manuscript sources and
figures, interactive numerical visualizations, Lean sources and
formalization reports, and the timestamp archive:

\noindent\url{https://github.com/zhangyuetong/smooth-tight-nonrigidity}
